\documentclass[12pt,a4paper]{article}
\usepackage[utf8]{inputenc}
\usepackage[T2A]{fontenc}
\usepackage[russian]{babel}
\usepackage{amsmath,amssymb}
\usepackage{array}
\usepackage{longtable}
\usepackage{booktabs}
\usepackage{geometry}
\usepackage{hyperref}
\geometry{margin=2.3cm}

\title{Углублённый аудит гипотезы единства для \texttt{TOE.pdf} и \texttt{UGSM}}
\author{}
\date{}

\begin{document}

\maketitle

\begin{abstract}
Этот документ делает следующий шаг после первичного моста между \texttt{TOE.pdf} и \texttt{ugsm\_dynamics\_audit-3.pdf}. Если в предыдущем аудите была проверена сама возможность читать обе конструкции как две тени одного спектрального источника, то здесь исследуется более сильный вопрос: устойчива ли эта картина при изменении окна фита, при деформации радиусов \(S^3\times S^1\) и при замене конкретного теплового ядра на более общее семейство спектральных функций. Цель документа --- не доказать полную эквивалентность теорий, а проверить, сохраняется ли гипотеза единства при более жёстком численном стрессе.
\end{abstract}

\section{Что именно проверялось глубже}

В углублённом цикле были проверены три вещи.
\begin{enumerate}
    \item Устойчивость коэффициентов теплового следа при смене окна малого параметра \(t\).
    \item Устойчивость геометрической интерпретации при изменении радиусов \(R_3\) и \(R_1\).
    \item Сохранение единого спектрального характера при переходе от одного теплового ядра к нескольким спектральным kernels.
\end{enumerate}

Идея проста: если TOE и UGSM действительно отражают одну большую сущность, то картина не должна разрушаться от малых деформаций численной процедуры и от естественных изменений геометрического фона.

\section{Метод}

Как и ранее, рассматривался скалярный лапласиан на
\[
M=S^3(R_3)\times S^1(R_1),
\]
с собственными значениями
\[
\lambda_{n,m}=\frac{n(n+2)}{R_3^2}+\frac{m^2}{R_1^2},
\qquad g_n=(n+1)^2.
\]

Численный след вычислялся по формуле
\[
K(t)=\sum_{n,m} g_n e^{-t\lambda_{n,m}},
\]
после чего аппроксимировалась величина
\[
\widetilde K(t)=(4\pi t)^2 K(t) \approx a_0+a_2 t+a_4 t^2.
\]

В отличие от первого аудита, теперь использовались:
\begin{itemize}
    \item более глубокое усечение спектра \(n_{\max}=100\), \(m_{\max}=100\);
    \item три окна фита: narrow, baseline и wide;
    \item пять пар радиусов \((R_3,R_1)\);
    \item несколько спектральных kernels: \(e^{-x}\), \(e^{-2x}\), \((1+x)^{-3}\), \((1+x)^{-4}\).
\end{itemize}

Результаты сохранены в файле \texttt{spectral\_unity\_deep\_results.json}, а сам расчёт выполнен в \texttt{spectral\_unity\_deep\_checks.py}.

\section{Теоретические ожидания для радиусного прогона}

Для выбранного скалярного теста естественные ожидаемые коэффициенты имеют вид
\begin{align}
 a_0^{exp} &= 4\pi^3 R_3^3 R_1,\\
 a_2^{exp} &= 4\pi^3 R_3 R_1.
\end{align}

Это означает следующее:
\begin{itemize}
    \item объёмный сектор должен масштабироваться как \(R_3^3R_1\);
    \item кривизновый сектор должен масштабироваться как \(R_3R_1\);
    \item систолический \(\pi\)-proxy должен сохраняться как \(\pi R_1\).
\end{itemize}

Если численный эксперимент действительно ловит общий геометрический источник, то именно такие законы масштабирования и должны проявиться.

\section{Проверка устойчивости по окну фита}

Для базовой геометрии \(R_3=R_1=1\) были взяты три окна:
\begin{itemize}
    \item narrow: \(t=0.02\) -- \(0.04\);
    \item baseline: \(t=0.03\) -- \(0.06\);
    \item wide: \(t=0.04\) -- \(0.08\).
\end{itemize}

\begin{center}
\begin{tabular}{lccc}
\toprule
Окно & $a_0$ relative error & $a_2$ relative error & $a_2/(4\pi)$ \\
\midrule
narrow & $4.17\times 10^{-6}$ & $4.45\times 10^{-4}$ & 9.8652145086 \\
baseline & $1.44\times 10^{-5}$ & $1.01\times 10^{-3}$ & 9.8596 \\
wide & $3.41\times 10^{-5}$ & $1.82\times 10^{-3}$ & 9.8516829207 \\
\bottomrule
\end{tabular}
\end{center}

\subsection{Вывод по окнам}

Картина остаётся устойчивой. Даже при расширении окна ошибки для \(a_0\) остаются на уровне \(10^{-5}\), а для \(a_2\) --- на уровне \(10^{-3}\). Это означает, что базовая интерпретация моста не является артефактом одного специально подобранного диапазона \(t\). Напротив, она сохраняется при разумной деформации процедуры фита.

\section{Радиусный прогон}

Следующий шаг состоял в проверке, сохраняются ли объёмный и кривизновый законы при изменении геометрии. Были исследованы пары радиусов
\[
(1.0,1.0),\ (1.2,1.0),\ (1.0,1.3),\ (1.5,0.8),\ (0.8,1.5).
\]

\begin{longtable}{p{2.2cm}p{2.2cm}p{2.5cm}p{2.5cm}p{2.3cm}p{2.3cm}}
\toprule
$R_3$ & $R_1$ & $a_0$ rel. err. & $a_2$ rel. err. & $a_2/(4\pi)$ fit & $a_2/(4\pi)$ expected \\
\midrule
\endhead
1.0 & 1.0 & $1.44\times 10^{-5}$ & $1.01\times 10^{-3}$ & 9.8596 & 9.8696 \\
1.2 & 1.0 & $4.76\times 10^{-6}$ & $4.84\times 10^{-4}$ & 11.8378 & 11.8435 \\
1.0 & 1.3 & $1.44\times 10^{-5}$ & $1.01\times 10^{-3}$ & 12.8175 & 12.8305 \\
1.5 & 0.8 & $1.24\times 10^{-6}$ & $1.97\times 10^{-4}$ & 11.8412 & 11.8435 \\
0.8 & 1.5 & $5.58\times 10^{-5}$ & $2.52\times 10^{-3}$ & 11.8137 & 11.8435 \\
\bottomrule
\end{longtable}

\subsection{Вывод по радиусам}

Главный результат здесь очень сильный: ожидаемое масштабирование сохраняется не только качественно, но и количественно. Ведущий коэффициент \(a_0\) воспроизводит закон \(R_3^3R_1\) с ошибками порядка \(10^{-6}\) -- \(10^{-5}\), а коэффициент \(a_2\) воспроизводит закон \(R_3R_1\) с ошибками порядка \(10^{-4}\) -- \(10^{-3}\).

Это особенно важно для гипотезы единства. Если бы связь TOE и UGSM была только числовой случайностью при \(R_3=R_1=1\), то она легко разрушилась бы при деформации радиусов. Но этого не происходит: структура \(4\pi^3\)- и \(\pi^2\)-секторов остаётся привязанной к геометрии именно так, как и должна для одного общего спектрального источника.

\section{Проверка на семействе спектральных kernels}

Чтобы выйти за пределы одного теплового ядра \(e^{-x}\), были проверены ещё три спектральных kernels:
\[
e^{-2x},\qquad (1+x)^{-3},\qquad (1+x)^{-4}.
\]

Для каждого kernel вычислялось спектральное действие
\[
S_f(\Lambda)=\mathrm{Tr}\, f(D^2/\Lambda^2)
\]
при \(\Lambda=6,8,10,12\), а затем рассматривалась объёмно-нормированная комбинация
\[
\frac{S_f(\Lambda)}{\Lambda^4 R_3^3R_1}.
\]

\begin{center}
\begin{tabular}{lcc}
\toprule
Kernel & Mean normalized value & Relative spread on $\Lambda=6\ldots 12$ \\
\midrule
$e^{-x}$ & 0.79736 & 2.09\% \\
$e^{-2x}$ & 0.20239 & 4.19\% \\
$(1+x)^{-3}$ & 0.39343 & 3.83\% \\
$(1+x)^{-4}$ & 0.13494 & 4.30\% \\
\bottomrule
\end{tabular}
\end{center}

\subsection{Вывод по kernels}

Этот тест показывает важную вещь: меняя kernel, мы не теряем общий геометрический контроль. Абсолютная нормировка, разумеется, зависит от выбора \(f\), но объёмно-нормированная структура остаётся стабильной и медленно дрейфует по мере роста \(\Lambda\), что соответствует ожиданию от одного и того же спектра, читаемого через разные функциональные окна.

Именно это и нужно для моста TOE--UGSM: не буквенная неизменность всех чисел, а устойчивость того факта, что различные спектральные чтения продолжают видеть один и тот же геометрический каркас.

\section{Что это даёт для гипотезы единства}

После углублённого аудита можно сформулировать несколько более сильных утверждений, чем в первом документе.

\begin{enumerate}
    \item Гипотеза общей тени устойчива к смене окна фита; следовательно, она не выглядит артефактом одной локальной настройки.
    \item Гипотеза общей тени устойчива к деформации радиусов; следовательно, совпадение не похоже на случайный эффект только в единичной геометрии.
    \item Гипотеза общей тени устойчива к смене spectral kernel; следовательно, мост между TOE и UGSM имеет смысл не только для одного теплового представления, но и для более общего семейства спектральных чтений.
\end{enumerate}

В совокупности это существенно усиливает тезис о том, что обе теории действительно могут быть двумя проекциями одной более глубокой сущности.

\section{Что всё ещё не закрыто}

Даже после этого усиления сохраняются ключевые открытые пункты:
\begin{itemize}
    \item весь анализ пока выполнен в скалярном секторе, а не для физически более сильного оператора Дирака-типа;
    \item не доказано, что полный корреляционный оператор TOE строго совпадает с операторным источником UGSM, а не только редуцируется к нему в минимальном секторе;
    \item линейный \(\pi\)-сектор всё ещё имеет статус геометрически сильной интерпретации, но не полного строгого вывода;
    \item мост пока подтверждён как стабильная исследовательская программа, а не как завершённая редукция двух корпусов к одной теории.
\end{itemize}

\section{Итоговый вердикт}

Углублённый аудит усиливает исходный документ заметно сильнее, чем просто ещё одна таблица чисел. Теперь можно сказать следующее:

\begin{quote}
Гипотеза о том, что TOE и UGSM являются двумя тенями одной более общей спектральной сущности, прошла не только первичную численную проверку, но и проверку на устойчивость. Она сохраняется при изменении окна фита, при деформации геометрических радиусов и при переходе к нескольким спектральным kernels. Это делает тезис об их едином источнике существенно более сильным и методологически более серьёзным.
\end{quote}

Самое аккуратное резюме таково: на текущем этапе мы всё ещё не доказали единство окончательно, но уже существенно сузили пространство альтернативы ``это просто числовая случайность''. После проведённого прогона гипотеза единого спектрального источника выглядит заметно более правдоподобной и более устойчивой, чем после первого базового аудита.

\end{document}